\chapter{Tracking the Evolution of Individual Resonances}\label{chap:resonances}

Prepared for submission to an academic journal.

Harlan Heilman, Thomas Ferron, A. T. M. Mahbub Alahe, Camille Bishop, and Brian A. Collins. In Prep. (2026)

This chapter reproduces a manuscript in preparation. The second question of \cref{chap:introduction} is examined here for Alq\textsubscript{3} and ZnPc films, building on the oscillator description of the optical tensor in \cref{chap:molecular-films} and on the ZnPc structure refined in \cref{chap:znpc}.

\section{Attributions}\label{sec:attributions-1}

HH and TJF conceived the study. BAC and TJF oversaw and directed all aspects of the study. HH performed DFT Calculations, analyzed RSoXS data, developed \& performed the analysis routines, and derived the Basis functions. TJF collected and processed RSoXS and NEXAFS measurements. A.T.M.M.A deposited the samples under the advising of CB.

\section{Introduction}\label{sec:introduction-1}

Resonant soft X-ray scattering (RSoXS) is a powerful technique for probing nano-scale structures within materials, that are otherwise limited by low contrast, by exploiting the large differences in optical response between materials near atomic absorption edges.\autocite{grabner2025resonant} The extreme energy dependence of the optical contrast provide a mechanism for contrast variation, enhancing sensitivity to otherwise inaccessible regions of buried nano-assemblies across scattering modalities - Scattering (RSoXS, P-RSoXS), Diffraction (CD-RSoXS, PE-RSoXS), and Reflectivity (RSoXR, P-RSoXR). Applied to complex material systems with functionalized applications, RSoXS has revealed structural mechanisms behind emergent performance driving their continued improvement from electronic devices such as solar cells and transistors to microelectronics and medical therapeutics. However, the driving force behind contrast enhancement, the fine structural spectral differences between materials near absorption edges, remains under explored with few studies performing comprehensive analyses on the chemical variations within interfaces in nano-assemblies.

This challenge rests jointly on experimental challenges, computational complexity and interpretive modeling. Most established, the optical response within buried scatterers is modeled from tabulated measurements performed on standard samples independently of the RSoXS experiment. Yet this approach fails for systems where standards cannot be measured such as volatile organic solvents, metals bearing novel oxidation states and molecules in unique conformations. Without material standards, the optical response must instead be handled on a per energy basis, the standard approach for RSoXR / P-RSoXR, a common structure is refined across energies, with the optical constants fit for every measured energy. From the measured optical constants, the task is to infer an underlying chemical composition or state within the scatterers. This approach is powerful but leads to a parameter explosion. Each newly measured energy increases the model complexity rather than improving the determination of the chemical state and interpreting the state itself is not straightforward. Because of this, the common solution is to reduce the number of measured energies, making the analysis tractable, but also removing the capability to resolve novel chemical states at the interfaces most critical to emergent function.

Fully characterizing novel chemical states within buried structures requires a framework that can produce the full optical spectra from a reduced set of parameters that can be refined along with the structure during analysis. The variational analysis presented by Kuzmenko offers insight into a solution that constrains the optical response using the Kramers-Kronig (KK) causality constraints. Expressing the optical response with KK consistent basis functions \(f_{i}\left( E;\overrightarrow{\theta_{i}} \right)\) i.e.

\begin{equation}
  n(E) = \sum_{i} f_{i}\left( E;\vec{\theta}_{i} \right)
  \label{eq:basis-sum}
\end{equation}

produces a set of auxiliary parameters \(\overrightarrow{\theta_{i}}\) that are unchanging with the number of measured energies\autocite{kuzmenko2005kramerskronig} Using a triangle basis, Stone et al. apply this method to directly uncover the optical constants from resonant reflectivity measurements collected over a full absorption edge - a powerful technique that has seen wide success.\autocite{stone2012kramers,macke2014magnetic,dipalo2024attosecond,corbae2026probing} Despite this success, the Kuzmenko-Stone basis falls prey to the same parameter explosion with energy sampling. Dense and time-consuming energetic sampling is necessary to acquire a high numerical precision across the absorption edge, and each measured energy gains its own basis function, so the parameter count grows with the energy sampling.

Here we parameterize the optical response with a basis of oscillators, directly mapping the underlying transition dipole moments responsible for the optical response. This choice collapses the number of basis functions needed to represent the spectral response, while enforcing Kramers-Kronig consistency at the same time as constraining the anisotropic response with the dipole moment. Using this approach, we demonstrate the capability to directly measure spectroscopic information from buried nano scale regions near the molecule-substrate interface across multiple molecular platforms. This new capability enables new quantitative analyses on the novel electronic, magnetic, and chemical properties within these extreme interfacial regimes within both static and dynamic processes. Most critically, these results establish a new convention for the chemical specificity of RSoXS techniques and reveal flaws in conventional approaches to RSoXS energetic sampling methods.

\section{Problems with Experimental Optical Models}\label{sec:problems-with-experimental-optical-models}

\begin{figure}[htp]
  \centering
  \includegraphics[width=\linewidth,height=0.75\textheight,keepaspectratio]{fig4.1_ML-DFT}
  \caption[Sample dependent NEXAFS for Alq\textsubscript{3}.]{Sample dependent NEXAFS for Alq\textsubscript{3}. (a) Imaginary refractive index component \(\beta(E)\) compared between samples of different deposition thicknesses presenting different spectral features along with ML clustered DFT calculations (gray). (b) KK calculated real refractive index component for the three different sources of \(\beta(E)\). (c) Peak amplitude difference map between the three models.}
  \label{fig:5.1}
\end{figure}



The dominant approach for constructing optical models for RSoXS originates from experimental measure of the NEXAFS spectroscopy on thin film standard samples of the material. The index of refraction \(n(E) = 1 - \delta(E) + i\beta(E)\) is found in pieces, \(\beta(E)\) is directly measured from the measured absorption spectra (\cref{fig:5.1}a), and \(\delta(E)\) is computed from the KK transformation (\cref{fig:5.1}b).\autocite{watts2014calculation,yan2013accurate} However, not all NEXAFS measurements probe the same absorptive response. In total electron yield mode (TEY) the incident X-ray beam excites only the top surface of the material, and electrons are collected through a conductive substrate, meaning the measured signal convolves surface effects with the material's own conductivity. For emergent materials bearing engineered optical and electronic properties, different preparations can drastically alter the measured NEXAFS spectra. Additionally, it is not always possible to form thin film standards, such is the case for organic solvents, metals bearing novel oxidation states, mixed phase interfaces, molecules in unstable conformations, and interface driven bonding hybridizations.

\Cref{fig:5.1} demonstrates this in thin vapor deposited films of Alq\textsubscript{3}, a common material platform used both as an electron transport and emissive layer in organic light-emitting devices, and in spintronic applications. Here we present two samples prepared under near identical conditions but with drastically different NEXAFS spectra. Notably within the \(\pi\) manifold where the relative peak intensities of all transitions vary strongly between samples, a response that is reflected similarly in \(\delta(E)\). In addition to this intensity only variation, a low energy transition at \qty{284}{eV} can be resolved in the thicker film but is entirely suppressed in the thin film NEXAFS. Machine learning clustered density functional theory (ML-DFT) calculations suggest the presence of a family of transitions at this energy capable of describing this feature. However, the presented spectra use the raw calculated transition energies with no static energy shifts, so the absolute energy scale is not well determined from DFT. A shift of \qty{+0.01}{eV} maximizes the alignment to the \qty{100}{nm} film, while a shift of \qty{0.8}{eV} would maximize alignment for the thinner \qty{20}{nm} sample, removing the possibility for the ML-DFT transitions to describe this \qty{284}{eV} transition, implying that this transition emerges as an effect of the thicker film and is not present in either a single molecule DFT calculation or within thinner films of the material. Regardless with the unshifted ML-DFT calculations, it is possible to refine the transition moments to a common set of transition clusters with only oscillator strength varying between the samples.

\section{The Spectral Basis}\label{sec:the-spectral-basis}

Optical responses in macroscopic molecular and polymer material systems emerge from their ensemble of molecules and polymers within the film. Assuming the packing characteristics of nearby molecules and polymer chains does not disrupt or alter the electronic states and excitation energies, then the individual molecular contributions sum linearly, weighted by their population representation within the ensemble. Generally, this can be expressed as the ensemble average over molecular configuration parameters \(\left\{ p_{i} \right\}\)

\begin{equation}
  n_{\mathrm{film}} = \left\langle n_{\mathrm{mol}} \right\rangle_{\{ p_{i} \}}.
  \label{eq:ensemble-average}
\end{equation}

Most commonly these handles molecular orientation degrees of freedom such as tilt angle or order parameters and packing characteristics such as density scaling. Multi component films are constructed similarly from an effective media approximation, but linear volume fraction weighted combinations is a good approximation for soft matter in the soft X-ray regime. Regardless this presents the key to transforming the molecular response into the film optical constants that are used to describe scattering.

The molecular resonant response can be decomposed into three regions, the non-resonant pre-edge, near edge fine structure, and the post edge. While the ionization phenomena that lead to the edge itself include nonlinear, multi electron and multi photon interactions, the edge description itself composeslinearly with the fine structure to describe the resonant response. Because of this, it can be separated from the fine structure

\begin{figure}
  \centering
  \includegraphics[width=\linewidth,height=0.75\textheight,keepaspectratio]{fig5-02_oscillator-basis}
  \caption[Gaussian spectral basis compared with a Lorentzian oscillator and triangle bases.]{(a) Spectral, Gaussian basis compared with Lorentzian damped driven oscillator and the sum of 26 triangle bases. (b) Sum over states approach constructing \(\beta(E)\) from the sum of Gaussian resonance building blocks, each originating from a refined ML-DFT cluster. (c) Derived \(\delta(E)\) computed directly from the sum of Dawson functions, overlayed the KK calculation of the results.}
  \label{fig:5.2}
\end{figure}


\begin{equation}
  n(E) = 1 - \left[ \delta_{\mathrm{edge}}(E) + i\beta_{\mathrm{edge}}(E) \right] + \sum_{i}\left[ \delta_{i}\left( E;\vec{\theta} \right) + i\beta_{i}\left( E;\vec{\theta} \right) \right],
  \label{eq:sos-index}
\end{equation}

where \(\beta_{i}\) is the absorption associated with a single resonant feature, and \(\delta_{i}\) its KK conjugate pair. This decomposition is commonly called a building block approach but is better described as a spectral decomposition as each resonant feature derives from the molecular quantum transitions. Each resonant feature represents an ensemble of quantum transitions from tightly bound atomic states to unoccupied molecular orbitals, so the line shape assumes the convolution of fundamental lifetime broadening and thermal broadening from the ensemble of transitions.\autocite{stohr2010nexafs} When lifetime broadening is dominant, resonant features follow the Lorentzian line shape derived from the textbook damped driven harmonic oscillator.\autocite{stohr2010nexafs,kuzyk2023nonlinear} Phenomenologically, this is rarely the case and ensemble broadening dominates resulting in a Gaussian line shape profile. While the real component \(\delta_{i}\) for the Lorentz oscillator can be derived from first principles, Gaussian bases require performing the KK transformation, resulting in the Dawson special function. This results in a functional pair

\begin{equation}
  \begin{gathered}
  \beta\left( E;E_{i},w_{i},A_{i} \right) = \frac{A_{i}}{w_{i}\sqrt{\pi}}\exp\left[ - \left( \frac{E - E_{i}}{w_{i}} \right)^{2} \right], \\
  \delta\left( E;E_{i},w_{i},A_{i} \right) = \frac{2A_{i}}{w_{i}\pi}D\left[ \frac{E - E_{i}}{w_{i}} \right],
  \end{gathered}
  \label{eq:gauss-dawson-pair}
\end{equation}

where \(w_{i} = \Gamma_{i}/(2\sqrt{\ln 2})\) for a given full width at half maximum \(\Gamma_{i}\), \(A_{i}\) the transition intensity, \(E_{i}\) the transition energy, and \(D\) the Dawson function.

This direct spectral decomposition presents directly interpretable parameters that map onto the individual quantum states associated with the resonant feature. The resonant energy immediately maps to the transition energy associated the quantum states, and the amplitude is directly proportional to the oscillator strength and TDM. Importantly, the transition intensity depends on the angle between the TDM orientation and the electric field polarization vector so that not only does this basis enforce the Kramers-Kronig relations, but it also enforces molecular, material, and dipole symmetries, producing the index of refraction tensor for free.

As a tensor the decomposition can now be augmented so that the peak amplitudes follow \(A_{i}\left( {\mu_{i},\vartheta}_{i},\varphi_{i} \right)\) where \(\vartheta_{i}\) and \(\varphi_{i}\) are the polar and azimuthal angles associated with the TDM orientation and \(\mu_{i}\) is the transition dipole moment amplitude (\cref{fig:5.2}e). This allows for the amplitude to be decomposed into an \(xx\), \(yy\), and \(zz\) component the full tensor. Projecting the transition dipole moment onto the lab frame through its polar and azimuthal angles gives \(A_{i,xx} = \mu_{i}^{2}\sin^{2}\vartheta_{i}\cos^{2}\varphi_{i}\), \(A_{i,yy} = \mu_{i}^{2}\sin^{2}\vartheta_{i}\sin^{2}\varphi_{i}\), and \(A_{i,zz} = \mu_{i}^{2}\cos^{2}\vartheta_{i}\). Once the dipole is described in the molecular reference frame each of the dipole orientations and their magnitude can be refined to match experimental results.

\section{Measuring Buried Electronic States}\label{sec:measuring-buried-electronic-states}

\begin{figure}[htp]
  \centering
  \includegraphics[width=\linewidth,height=0.75\textheight,keepaspectratio]{fig5-03_alq3-fit}
  \caption[Alq\textsubscript{3} reflectivity fit validating the resolved ML-DFT cluster model.]{Alq\textsubscript{3} reflectivity fit validating the resolved ML-DFT cluster model. (a) RSoXS reflectivity in s- and p-polarization at ten resonant energies (offset for clarity), with the resolved sum-over-states (SOS) model (solid) and the model fixed to the \qty{20}{nm}-film NEXAFS optical constant (dashed). (b) Expanded low-q region showing the SOS model tracking fringe positions and amplitudes that the \qty{20}{nm}-NEXAFS model misses. (c) Depth-resolved \(\beta(E)\) (top) and \(\delta(E)\) (bottom) recovered by the SOS model, colored by depth from the free surface to the substrate. (d) Density profile recovered by the SOS model (solid) compared with the \qty{20}{nm}-NEXAFS model (dashed), which is pushed to its fit bounds at the surface and in the bulk.}
  \label{fig:5.3}
\end{figure}

The Alq\textsubscript{3} optical model is defined by first refining the ML-DFT clusters to NEXAFS collected on the same sample (\cref{fig:5.1}). Peak amplitudes are then refined as a function of depth with widths and transition energies held fixed. The material structure is modeled with a single slab of Alq\textsubscript{3} with its density following a continuous exponential transition from the interface induced density at the surface and the free surface and substrate into the bulk of the film, with surface roughness modeled in the Névot-Croce sense.\autocite{heilman2026sub,nevot1980caracterisation} This approach uses pinning densities \(\rho_{Si}\ \&\ \rho_{Surf}\) with respective exponential characteristic thicknesses \(\tau_{Si}\ \&\ \tau_{Surf}\) along with the bulk properties \(\rho_{bulk}\), total thickness \(t\) and surface roughness \(\sigma_{r}\). The Alq\textsubscript{3} peak heights then vary over this same functional depth profile with bulk and interfacial pinned amplitudes for each refined peak, tying the electronic deviations within Alq\textsubscript{3} molecules to the packingcharacteristics set by the interfacial surface potentials, a feasible mechanism for measured depth dependent optical responses.



\Cref{fig:5.3} showcases the model fit results on RSoXS reflectivity measured at ALS beamline 11.0.1.2. High \(q\)-range measurements at 284.5 and \qty{287.2}{eV} (\cref{fig:5.3}a \& vertical lines in \cref{fig:5.3}c) maximize spatial resolution with the aim to capture small and buried structures and are captured at energies with low resonant absorption \(\beta\) and with opposite signs of \(\delta\). The shorter \(q\)-range measurements of \cref{fig:5.3}b aim to better capture the general spectral shape (shaded \cref{fig:5.3}c). Two optical models are compared, the static ML-DFT NEXAFS refined optical constants (dashed throughout) against the SOS depth dependent response (solid traces). The sum over states (SOS) model, uses 28 free parameters describing the evolution of five resonant transition amplitudes along the depth of the film, compared with the static model derived from NEXAFS collected on the same film. The resonance refinements improve the reduced \(\chi^{2}\) best fit statistic \(\chi_{\nu}^{2}\) from 63.1 to 12.9 and improves the BIC by 1527 (\cref{tab:5.1}), despite having more than twice as many free parameters. This is not the result of a single poorly described energy, the weighted RMSE of the NEXAFS-fixed model exceeds the SOS model at every one of the ten measured energies, roughly by a factor of two to three throughout the \(\pi^{*}\) manifold.

\begin{table}
  \centering
  \caption{Fit quality comparison between the SOS and static NEXAFS models.}\label{tab:5.1}
  \begin{tabular}{l S[table-format=2.0] S[table-format=2.1] S[table-format=1.2] S[table-format=4.0]}
    \toprule
    Model & {\(k\)} & {\(\chi^{2}_{\nu}\)} & {Weighted RMSE} & {\(\Delta\)BIC} \\
    \midrule
    SOS (5 peaks) & 28 & 12.9 & 3.55 & 0 \\
    TEY NEXAFS & 13 & 63.1 & 7.89 & 1527 \\
    \bottomrule
  \end{tabular}
\end{table}

The refined model reproduces both the measured \(\beta(E)\ \)and its KK-derived \(\delta(E)\) directly as a function of depth within the film (\cref{fig:5.3}c) derived as an evolution of resonant intensities along the depth of the film. Generally, these results describe an increase in absorption uniformly across all peaks as a function of depth, with the strongest absorption at the substrate interface. The static optical approach attempts and fails to capture this behavior using a strong density gradient pushing the density at the substrate to 2.5 times what its stoichiometry and packing geometry allow. The strongest evolution arises with the third resonance, nearly fully suppressed right at the vacuum surface, before growing substantially into the bulk and at the silicon interface. The origin of this evolution is not well understood and may arise from any of, hydrocarbon contamination at the surface so that this interfacial resonant state is a model hallucination, molecular conformation shifts that suppress the transition rates near the surface while enhancing them near the substrate, the surface potentials themselves altering the electronic structure leading to hybridization, or some other optical or electronic property of Alq\textsubscript{3} not characterized in these films. The challenge is that the present data only directly measures the optical constants at the energies probed with the optical constants at other energies inferred based on the dispersive properties of each resonant transition, with most of this out of measurement sensitivity originating from \(\delta(E)\). To properly investigate the source of the depth dependent optical evolution; better sampling strategies are needed to improve the reliability of these measured optical constants.

\section{Tracing Dipole Evolution with Depth}\label{sec:tracing-dipole-evolution-with-depth}

\begin{figure}[htp]
  \centering
  \includegraphics[width=\linewidth,height=0.75\textheight,keepaspectratio]{fig5-04_znpc-tdm-depth}
  \caption[Depth-resolved transition dipole moment orientation in ZnPc.]{Depth-resolved transition dipole moment orientation in ZnPc. (a) P-RSoXR reflectivity (zoomed ×3, offset for clarity) at five resonant energies in s- and p-polarization, with the resolved sum-over-states (SOS) fit (black) and a NEXAFS-fixed fit (gray). (b) Molecular xx and (c) molecular zz tensor components of \(\beta(E)\) (top) and \(\delta(E)\) (bottom) recovered by the SOS model at each sampled depth (colorbar). (d) Density profile \(\rho(z)\) (solid, bottom axis) and the resolved TDM tilt angle \(\gamma(z)\) (dashed, top axis), with the free-surface and SiO\textsubscript{2}/Si substrate regions shaded.}
  \label{fig:5.4}
\end{figure}

Next, we demonstrate the approach for tracking individual dipole orientations along the depth of the film. Zinc phthalocyanine (ZnPc) is a well-studied molecular platform that forms novel interfacial molecular states when deposited on silicon.\autocite{heilman2026sub} Additionally, the ML-DFT tensor of ZnPc reproduces angle-resolved NEXAFS of these films with an average error of 2.7 percent and, with a single tilt and density at each depth, describes their polarized reflectivity at every measured energy (\cref{chap:znpc}).\autocite{murcia2026quantitative,heilman2026sub} Moreover, PVD can easily produce highly anisotropic films, contrary to Alq\textsubscript{3} making it a natural platform to evaluate capability for this technique to track the TDM orientation of individual dipoles along with the molecular orientation as a function of depth within thin films.



\Cref{fig:5.4}a shows polarized reflectivity captured from \qty{20}{nm} thin films of ZnPc deposited on bare silicon (the film of \cref{chap:znpc}\autocite{heilman2026sub}). Following the similar convention described above, a long q-range measurement is performed at \qty{283.7}{eV} with low absorption, yet this time directly probing the energies of the four most prominent resonances within the \(\pi\) manifold (ticks, \cref{fig:5.4}). The structural model follows the same single slab continuous profile as previously published, with the orientation and density at the surface and substrate interfaces set by their respective surface potentials, with orientation and density relaxing into the bulk. The optics are then extended by also fitting the TDM amplitude of each of the four primary resonant transitions, with depth dependence set by the total thickness \(t\) and same characteristic interfacial thicknesses \(\tau_{Si}\ \&\ \tau_{Surf}\). The molecular optical tensor is highly uniaxial, so each TDM is symmetrized fitting only the polar cone angle \(\vartheta\ \)for each of these four resonances. In the fits reported here each cone angle is held constant with depth, so peak amplitudes vary over the film while the tensor projection of each resonance changes with depth only through the molecular tilt profile of \cref{chap:znpc}. As with Alq\textsubscript{3}, the resolved model is checked against a conventional NEXAFS-fixed alternative. \Cref{tab:5.2} compares the two fits to the \(N = 745\) point ZnPc reflectivity dataset shown above (\qty{283.7}{eV} and four resonant energies, s and p polarization, measured counting-statistics errors, fit on log-reflectivity residuals). model again fits markedly better, χ²ν = 3.63 against 10.33 and ΔBIC = 690 in its favor, but the per-energy residuals show a different pattern than Alq\textsubscript{3}: the two models agree closely at the lower three energies and diverge sharply only at 287.0 and \qty{289.0}{eV}, where the NEXAFS-fixed weighted RMSE more than doubles. These are the energies carrying the strongest TDM-orientation sensitivity in this molecule, so the conventional model is adequate where orientation contrast is weak and fails specifically where it is strong, rather than failing uniformly.

\begin{table}
  \centering
  \caption{Fit quality comparison between the SOS and static NEXAFS models.}\label{tab:5.2}
  \begin{tabular}{l S[table-format=2.0] S[table-format=2.2] S[table-format=1.2] S[table-format=3.0]}
    \toprule
    Model & {\(k\)} & {\(\chi^{2}_{\nu}\)} & {Weighted RMSE} & {\(\Delta\)BIC} \\
    \midrule
    SOS (4 peaks) & 37 & 3.63 & 1.86 & 0 \\
    NEXAFS & 21 & 10.33 & 3.17 & 690 \\
    \bottomrule
  \end{tabular}
\end{table}

Four resonances were opened in the fit, numbered 1 to 4 in order of energy at 284.2, 285.1, 287.1, and \qty{288.9}{eV}, with the molecular tilt and density held at the \cref{chap:znpc} gradient values (\ang{58.7}, \ang{69.6}, and \ang{38.4} and 2.05, 1.64, and \qty{0.87}{g.cm^{-3}} at the vacuum surface, in the bulk, and at the SiO\(_{2}\) interface). Dividing each fitted height by the local density and by its NEXAFS-refined amplitude gives a per-molecule oscillator strength relative to NEXAFS. The two lowest \(\pi^{*}\) resonances carry the largest oscillator strength and move in opposite directions. At the vacuum surface resonance 1 falls to 0.30 of its NEXAFS value while resonance 2 rises to 1.45, and at the SiO\(_{2}\) interface resonance 1 rises to 1.65 while resonance 2 falls to 0.67, so between the two boundaries resonance 1 grows 5.5-fold and resonance 2 halves. Their summed strength stays within 4 percent of the NEXAFS value at both boundaries (1.04 at vacuum, 1.02 at SiO\(_{2}\)) and rises to 1.19 in the bulk, an excess reported with the caveat that density and tilt are fixed at the \cref{chap:znpc} values. Resonance 3 is unchanged at the surface and in the bulk (0.95) and falls to 0.25 at the interface, and resonance 4 changes most in relative terms, reaching 2.90 at the interface where the \(\pi^{*}\) manifold meets the \(\sigma^{*}\) onset. Each opened resonance carries one cone angle, constant with depth. The \(\pi^{*}\) angles stay within \ang{9} of the ring normal (resonance 1 from \ang{3.0} in the NEXAFS refinement to \ang{0.0}, resonance 2 from \ang{0.9} to \ang{8.9}), while resonance 3 moves from \ang{28.5} to \ang{24.6} and resonance 4 from \ang{26.0} to \ang{12.3}, so the largest reorientation relative to NEXAFS also occurs at the \(\pi^{*}\) to \(\sigma^{*}\) boundary.

In the ML-DFT clusters, resonance 1 is dominated by 1s excitation of carbon site C2 into the lowest unoccupied orbital (MO 117), and resonance 2 by excitation of site C3 into a higher \(\pi^{*}\) orbital (MO 119), both benzo-ring carbons, while resonance 4 arises from site C1, the pyrrole carbon bound to nitrogen, whose 1s ionization potential lies \qty{1.5}{eV} above the benzo carbons (site labels as in Supplemental Material). Because the structural model already fixes the local density, we surmise that the anti-correlated change of resonance 1 and resonance 2 at the two interfaces, with a nearly conserved sum, is a redistribution of spectral weight between the two lowest unoccupied states rather than a change in the number of absorbing molecules. Three mechanisms could produce it, each with a different signature. Hybridization or charge transfer at the SiO\(_{2}\) interface would perturb the LUMO most where molecules contact the oxide, confining the change to the first molecular layers and shifting the transition energy as well as its strength. Packing-dependent \(\pi\) overlap tied to the tilt gradient would scale the reweighting with the local tilt \(\gamma(z)\), so the same ratio of resonance 1 to resonance 2 should recur wherever the tilt is the same. Disorder at the free surface would broaden the resonances and lower the anisotropy of every transition rather than exchange weight between two. A refit that frees each cone angle separately at the surface, in the bulk, and at the interface, together with a fixed-angle energy scan from 283 to \qty{290}{eV} at several incident angles, would separate these, since only the first predicts an energy shift and only the second ties the reweighting to \(\gamma(z)\).

These results also open a route to refining the ML-DFT clusters themselves, using the clusters as computed by the clustering routine as priors for the spectral analysis. The energetic uncertainties of the present datasets are too large for this, and fixed-angle reflectivity across the edge at one or more incident angles, measured with this aim, may suffice.

\section{Energetic Sampling Efficiency}\label{sec:energetic-sampling-efficiency}

\begin{figure}[htp]
  \centering
  \includegraphics[width=\linewidth,height=0.75\textheight,keepaspectratio]{fig5-05_energy-sampling}
  \caption[Monte-Carlo sampling of resonant transitions for polystyrene (C\textsubscript{8}H\textsubscript{8}).]{Monte-Carlo sampling of resonant transitions for polystyrene (C\textsubscript{8}H\textsubscript{8}). Generated \(\beta(E)\) (a) and conjugate \(\delta(E)\) (b) traces over flat priors. Circles indicate fixed energy measurement that perfectly determine both δ and β at a given energy though measurement simulating RSoXS analysis at that energy. Remaining traces are NEXAFS profiles that agree at the measured energies with the ground truth in bold.}
  \label{fig:5.5}
\end{figure}



One natural consequence of this analysis is a framework for generating feasible, KK-consistent optical constants and tensors. The number of peaks, along with their heights, widths, and transition energies, can be sampled over a realistic parameter domain, with the optical response composed from their sum. This approach may be useful for generating training data in machine learning, and it also presents a natural way to explore sampling strategies and measurement uncertainty within resonant techniques, as is done in \cref{fig:5.5}. A synthetic model of ten carbon K-edge transitions on a fixed bare-atom continuum (constructed for a stoichiometry of C\textsubscript{8}H\textsubscript{8}) serves as the ground truth. Transition energies are drawn from Gaussian priors, and widths and heights from log-normal priors, centered on the ground truth. Each sampling strategy aims to measure \(\beta\) and \(\delta\) at its chosen energies within a Gaussian error of \(\sigma\  = \ 10⁻⁴\), simulating a sparse RSoXS measurement that precisely extracts optical constants at only a few energies. The family of \(\beta(E)\) and \(\delta(E)\) curves consistent with those measurements is then drawn from the Bayesian posterior by Markov chain Monte Carlo. All curves agree with the measurement error but diverge between them, visualizing the uncertainty a given energy sampling leaves unresolved. We quantify this with the fractional reduction, relative to the prior, of the variance in the complex index \(n\  = \ 1\  - \ \delta\  + \ i\beta\) integrated over 283 to \qty{290}{eV}.

Three of the four strategies measure four energies. The first places them on the leading resonance centers (\cref{fig:5.5}a). The second draws them uniformly at random (\cref{fig:5.5}b). The third selects them greedily with a Gaussian process whose kernel is the prior covariance of \(\beta\) and \(\delta,\) with each energy chosen to maximize the expected drop in variance given the earlier picks (\cref{fig:5.5}c). The fourth is a continuous energy scan in \qty{0.1}{eV} steps across the same window (\cref{fig:5.5}d) but using a wider uncertainty to simulate an energy scan such as a fixed angle reflectivity measurement. Measuring on the resonances reduces the variance by 78.0\%, but above \qty{286}{eV} the posterior traces fan out freely in both \(\beta\) and \(\delta\). The random set shown reduces it by 69.7\% and leaves the 285 to \qty{286.5}{eV} region unconstrained. This seed is favorable. Across 200 random energy sets, a linear-response estimate gives 53 ± 19\%. The Gaussian-process design moves one point onto the flank between resonances at \qty{285.2}{eV} and another out to \qty{287.2}{eV}, raising the reduction to 85.6\%. Even so, four energies supply eight observables against 30 transition parameters. Under the same linear-response estimate, a tenfold improvement in precision leaves the on-resonance reduction unchanged near 83\%. The equal-dose energy scan reaches 98.7\% despite its higher per-point noise.

For sparse RSoXS measurements, the limit is therefore set by the number of energies and not the precision at a single energy. Even exact knowledge of β and δ at four energies leaves the spectrum between them undetermined, with an optimized choice of energies recovering only part of that gap. We surmise that a model-based analysis of sparse-energy RSoXS should therefore be anchored to a continuous energy dependent scan even if each measurement over the scan has higher measurement noise and resolves \(\delta\) and \(\beta\) to a lower precision. Short of that, the same sampling framework can be used to choose measurement energies before a beamtime, extending traditional ideas selecting energies based on contrast.

\section{Discussion}\label{sec:discussion}

The Alq\textsubscript{3} and ZnPc case studies demonstrate the same effect from two different directions. In Alq\textsubscript{3}, replacing a single borrowed NEXAFS optical constant with independently refined resonant-feature heights removes a degeneracy between optical model and film density that the simpler model had been resolving by pushing density to an unphysical value. In ZnPc, replacing a single borrowed tensor with per-resonance cone angles and depth-resolved amplitudes reveals that different transitions reorient to different degrees relative to their NEXAFS-refined values and that the two lowest π* resonances exchange spectral weight between the two interfaces, a distinction the rigid model cannot express at all. Neither result follows from overfitting and both follow from adding optical degrees of freedom selectively to the independent resonances that poorly describe the measured signal. Tracking these resonances through the depth of the film then exposes physics, such as changes in packing or hybridization near an interface, that no single reference spectrum can capture.

Nothing in this spectral-decomposition framework is specific to Alq\textsubscript{3}, ZnPc, or small molecules, and the same analysis applies to any film, molecular, polymeric, or otherwise. All that is needed is a prior guess of the number of peaks needed to describe the optical response, and measurements dense enough in energy to constrain the optics across an absorption edge. By comparing model performance metrics including performance against non-fit validation data between best fit models constructed under different priors, it is possible to construct the optical constants directly without the need for the predefined resonant peaks that are taken here from the ML-DFT calculations. However, without these calculations, additional analysis on the final extracted resonant spectra is needed to understand the source of the electronic evolution throughout a film.

\section{Conclusion}\label{sec:conclusion}

We have presented a spectral decomposition of the optical response into Kramers-Kronig consistent oscillators tied to ML-DFT transition clusters, which refines the optical constants of buried regions together with the film structure. Because the parameter count is set by the number of resonances and not by the number of measured energies, resonant reflectivity becomes a direct probe of electronic states that no reference measurement reaches, extending its role beyond contrast variation. In Alq\textsubscript{3} the depth-dependent amplitudes remove a degeneracy between the optical model and the film density that the NEXAFS-fixed model could resolve only by driving the substrate density to an unphysical value. In ZnPc the per-resonance amplitudes and cone angles markedly improve the fit and show the two lowest π* resonances exchanging oscillator strength between the free surface and the oxide interface while their sum is nearly conserved. The sampling study shows that for sparse-energy measurements the number and placement of energies, and not the precision at each, limit the recovered spectrum, so continuous energy scans should anchor future analyses. Each refined resonance traces back to a family of electronic transitions, so the origin of its evolution with depth can be interrogated directly from the optical model, opening a route to depth-resolved spectroscopy of buried molecular layers.

\section{Methods}\label{sec:methods-1}

\textbf{Density functional theory:} A time independent transition potential calculation using the StoBe library was performed on a single molecule of ZnPc and Alq\textsubscript{3}. A HCH approach sets the manifold of transition energies, along with each TDM and oscillator strength, while the \(\Delta KS\ \)method is used to fix the absolute energy associated with each transition. For details see our previous work.\autocite{murcia2026quantitative} The ZnPc molecular geometry was extracted from a crystallographic structure file (.cif)\autocite{mai2020data} while Alq\textsubscript{3} used a thermally relaxed structure computed with UFF minimization.

Clustering was performed using the dft-learn implementation of the clustering routine from Murcia et.al. using energetic overlap to identify transitions.\autocite{murcia2026quantitative,victormurcia2025wsu} Oscillator strength threshold is determined using an elbow method. Overlap percentage is determined using a Bayesian Information analysis to select a percentage that produces the smallest set of clusters that best describe the original DFT calculation.

\textbf{Film growth and pre-characterization:} ZnPc samples were grown by physical vapor deposition under the same conditions described in prior work.\autocite{heilman2026sub} Small molecule ZnPc (C32H16N8Zn sourced from Sigma Aldrich) with purity \(\geq 99.5\%\) was further purified in the physical vapor deposition (PVD) chamber (base pressure \(< 1 \times 10^{- 7}\) Torr) by gradually heating, approaching ZnPc sublimation temperature over approximately an hour until no further off gassing of contaminants were detected by the quartz crystal monitor (QCM). Thin films (\({\sim}\)\qty{20}{nm} measured by QCM) of ZnPc were deposited at \(.2\ \mathring{\mathrm{A}}/s\) on high conductivity polished Si wafers (University Wafers). Wafers were cleaved \qty{1}{cm} x \qty{3}{cm} and cleaned via sonication in deionized water, acetone and isopropyl alcohols before a 30 min UV ozone bake. ZnPc was deposited onto these wafers. Alq\textsubscript{3} samples were grown by physical vapor deposition under similar conditions described in prior work.\autocite{ferron2023composition} Sublimation grade (99.995\% trace metals basis) Alq\textsubscript{3} was purchased from Sigma-Aldrich and deposited onto silicon substrates. Wafers were cleaved to and precleaned under the same conditions. Deposition was performed with the same base pressure and with thickness monitored to achieve \qty{24}{nm} of film growth.

\textbf{NEXAFS Collection:} ZnPc NEXAFS measurement methods are described in prior work with data taken from X-ray Atlas.\autocite{heilman2026x} Alq\textsubscript{3} NEXAFS on the \qty{200}{nm} sample are described in prior work with data taken from X-ray Atlas.\autocite{ferron2023composition} NEXAFS on the Alq\textsubscript{3} films studied independently though resonant reflectivity was performed in TEY mode at ALS beamline 11.0.1.2. Incident angles 30, 55 and 70 degrees demonstrate no dichroism within the film, consistent with DFT calculations.

\textbf{RSoXS Collection:} Polarized resonant soft x-ray reflectivity was measured at beamline 11.0.1.2 of the Advanced Light Source, Lawrence Berkeley National Laboratory, following the procedure established by Ferron et al.\autocite{ferron2023composition} Samples are transferred from inert storage to the end station where specular reflectance is captured on the CCD area detector with thermoelectric cooling. The sample surface is aligned to the goniometer axis of rotation by an iterative process that refines the sample rotation and height. Incident flux was attenuated using the scatter slits together with the higher order suppressor to avoid detector damage. Because the total reflectance varies by more than the detectors dynamic range, reflectivity is captured in overlapping \(\mathbf{q}\) subsets with suppressor settings and exposure time raised at each boundary once reflectivity reaches the noise floor. Data is collected first in s polarization and then in p polarization beginning at \qty{250}{eV}, below the carbon K edge, and stepping upward through the resonance. Upon polarization change, optics were allowed to re-equilibrate to the new heat load for 30 min before alignment was recalibrated. Photon energy was calibrated against the transmission NEXAFS of a polystyrene standard as well as electron yield on highly oriented pyrolytic graphite at the beginning and at the end of the shift with no discernable shift, bounding the energy drift over the experiment to the width of the primary polystyrene feature.

\textbf{Reflectivity modeling and fitting.} Reflectivity was computed with the 4 × 4 transfer matrix of refloxide and fit with refnx objectives.\autocite{alsrsoxsrefloxide2026als,nelson2019textless} For Alq\textsubscript{3} the film is a single slab whose density relaxes exponentially from pinned surface and substrate values into the bulk, with Névot-Croce roughness, and the fit is made on \(Rq^{4}\). For ZnPc the structure of \cref{chap:znpc} (tilt and density gradients, oxide, and substrate) is held fixed, and the fit frees the per-energy s and p scale factors and angle offsets, a global energy offset (refined to \(- 0.09\) eV), the heights of the opened resonances at the vacuum surface, in the bulk, and at the substrate, each bounded within 0.25 to 4 times its seed and non-negative at every depth, and one cone angle per opened resonance. Resonances were opened where they lie nearest the measured resonant energies. The data are the published ZnPc reflectivity at 283.7, 284.2, 285.1, 287.0, and \qty{289.0}{eV} in s and p polarization with measured uncertainties, fit on \(\log R\) by local least squares from the published graded fit, without a global search. The sampling study used Gaussian priors on transition energies and log-normal priors on widths and heights, a measurement error \(\sigma = 10^{- 4}\) on \(\beta\) and \(\delta\), posterior draws by MCMC, and a linear-response estimate over 200 random energy sets.
